% from main.tex
%%%%%%%%%%  scan idx 521  |  folio 508  %%%%%%%%%%
\textbf{Proposition 2.2.} Avec les notations et les hypothèses de 2.1, soit $f:P\to S$ le morphisme structural de $P$. Alors l'homomorphisme :
\[
f^*:K_\bullet(S)\longrightarrow K_\bullet(P)
\]
est surjectif, et son noyau est le sous-groupe de $K_\bullet(S)$ engendré par les sous-groupes $(\eta(m)-1)K_\bullet(S)$ pour $m\in M$, ou, ce qui revient au même, pour $m$ parcourant un système de générateurs de $M$.

On suppose d'abord $r=1$, i.e.\ $T=\mathbb G_{m,S}$. Alors $M=\mathbb Z$, et on a pour tout $m\in M$ :
\[
A_m\simeq A_1^{\otimes m}.
\]
Par suite, si on pose $A_1=L$, on a $A=\bigoplus_{n\in\mathbb Z}L^{\otimes n}$, ce qui montre que $P$ s'identifie au complémentaire de la section nulle dans le fibré $X=\mathbb V(L)$. Considérant $S$ comme un sous-schéma fermé de $X$ par la section nulle, on a d'après 1.1 la suite exacte :
\[
K_\bullet(S)\xrightarrow{g_*}K_\bullet(X)\longrightarrow K_\bullet(P)\longrightarrow 0,
\]
où $g:S\to X$ est la section nulle de $X$. Par ailleurs, $f^*$ se factorise en \(K_\bullet(S)\xrightarrow{f'^*}K_\bullet(X)\longrightarrow K_\bullet(P)\) où $f'$ est le morphisme structural de $X$. D'après 1.6, $f'^*$ est un isomorphisme, inverse de $g^*$. On a donc la suite exacte :
\[
K_\bullet(S)\xrightarrow{g^*\circ g_*}K_\bullet(S)\xrightarrow{f^*}K_\bullet(P)\longrightarrow 0,
\]
d'où l'on tire d'après 1.8 :
\[
K_\bullet(P)\simeq K_\bullet(S)/\lambda_{-1}(L)K_\bullet(S)=K_\bullet(S)/(1-L)K_\bullet(S),
\]
où $L$ est la classe dans $K^\circ(S)$ de $L$. Comme pour tout $n\in\mathbb Z$, $L^n-1$ est multiple de $1-L$, ceci prouve 2.2 dans le cas où $P$ est de rang $1$.

